\documentclass[11pt]{article}
\usepackage[a4paper,margin=1in]{geometry}
\usepackage{amsmath,amssymb}
\title{Sharp Summary: Campbell et al. (2022) -- Continuous-Time Discrete Diffusion}
\author{}
\date{}
\begin{document}
\maketitle

\section*{Problem and contribution}
This paper gives a full CTMC formulation of discrete diffusion: forward noising rates, reverse-time rates, continuous-time ELBO, and scalable sampling (tau-leaping + predictor-corrector).

\section*{Forward and reverse CTMC}
Forward infinitesimal transition:
\[
q_{t\mid t-\Delta t}(x\mid \tilde x)=\delta_{x,\tilde x}+R_t(\tilde x,x)\Delta t+o(\Delta t).
\]
Reverse-time CTMC is defined by
\[
q_{t\mid t+\Delta t}(\tilde x\mid x)=\delta_{\tilde x,x}+\hat R_t(x,\tilde x)\Delta t+o(\Delta t).
\]
For $x\neq \tilde x$:
\[
\hat R_t(x,\tilde x)
=R_t(\tilde x,x)
\sum_{x_0}
\frac{q_{t\mid 0}(\tilde x\mid x_0)}{q_{t\mid 0}(x\mid x_0)}
q_{0\mid t}(x_0\mid x),
\]
with $\hat R_t(x,x)=-\sum_{x'\neq x}\hat R_t(x,x')$.
This is the continuous-time analogue of Bayes inversion.

\section*{Learning objective}
Because $q_{0\mid t}$ is intractable, replace it by denoiser $p^\theta_{0\mid t}$ to get $\hat R_t^\theta$.
Then minimize the CT negative ELBO
\[
\mathcal L_{\mathrm{CT}}(\theta)
= T\,\mathbb E_{t\sim\mathcal U(0,T),\,q_t(x)r_t(\tilde x\mid x)}
\Big[
\sum_{x'\neq x}\hat R_t^\theta(x,x')
-\mathcal Z^t(x)\log \hat R_t^\theta(\tilde x,x)
\Big]+C.
\]
This is the continuous-time counterpart of discrete diffusion ELBO training.

\section*{Sampling theory and algorithmics}
Sampling uses tau-leaping, approximating jump counts by Poisson variables with means proportional to $\tau\hat R_t^\theta$.
Predictor-corrector adds a correction rate
\[
R_t^c = R_t + \hat R_t,
\]
whose stationary law at fixed $t$ is $q_t$.
The paper also provides an error decomposition combining reverse-rate approximation error and tau-leaping discretization error.

\section*{Relevance to this repo}
\textbf{Direct and high-value.}
It is the cleanest CTMC formalization of ``learn reverse dynamics from denoising posteriors'', and it informs how step-3 learning of $Q^t$ scales beyond toy finite-state chains.

\end{document}
